\begin{definitionbox}{Definition (Identity)}
\begin{definition}[Identity]
\label{def:identity-relation}
Two expressions $x$ and $y$ are \emph{identical} when they denote the same
object.
\end{definition}
\end{definitionbox}

\begin{remark*}[Standard quantified statement]
\[
x=y \Longleftrightarrow x \text{ and } y \text{ denote the same object}.
\]
\end{remark*}

\begin{remark*}[Interpretation]
Identity is sameness, not similarity. It is the relation used when no
distinction remains between two denotations.
\end{remark*}

\begin{dependencies}
\begin{itemize}
  \item \hyperref[def:relation]{Relation}
\end{itemize}
\end{dependencies}

\begin{axiombox}{Axiom (Leibniz's Law)}
\begin{axiom}[Leibniz's Law]
\label{ax:leibniz-law}
Identical objects have exactly the same properties.
\[
\forall x\,\forall y\,
\Bigl(x=y \Longrightarrow \forall P\,\bigl(P(x)\Longleftrightarrow P(y)\bigr)\Bigr).
\]
\end{axiom}
\end{axiombox}

\begin{remark*}[Standard quantified statement]
\[
\operatorname{Indiscernible}(x,y)
\Longleftrightarrow
\forall P\,\bigl(P(x)\Longleftrightarrow P(y)\bigr).
\]
\end{remark*}

\begin{remark*}[Interpretation]
Leibniz's law says that equality permits no observable mathematical
difference. Any property true of one equal object is true of the other.
\end{remark*}

\NoLocalDependencies

\begin{definitionbox}{Definition (Equality on a Domain)}
\begin{definition}[Equality on a Domain]
\label{def:equality-relation}
Let $A$ be a domain. \emph{Equality on $A$} is the relation, written $=$,
that holds exactly when two terms or constructions denote the same element of
$A$.
\end{definition}
\end{definitionbox}

\begin{remark*}[Standard quantified statement]
\[
\operatorname{Equality}_A(=)
\Longleftrightarrow
\forall x,y \in A,\,
\bigl(x=y \Longleftrightarrow x \text{ and } y \text{ denote the same element of } A\bigr).
\]
\end{remark*}

\begin{remark*}[Interpretation]
Equality is the identity relation restricted to the domain currently under
discussion. Later equivalence relations may identify objects for a purpose,
but equality records literal sameness.
\end{remark*}

\begin{dependencies}
\begin{itemize}
  \item \hyperref[def:relation]{Relation}
\end{itemize}
\end{dependencies}

\begin{definitionbox}{Definition (Definitional and Propositional Equality)}
\begin{definition}[Definitional and Propositional Equality]
\label{def:definitional-propositional-equality}
\emph{Definitional equality} records an introduced meaning, written with
$:=$. \emph{Propositional equality} is a statement of sameness inside the
theory, written with $=$.
\end{definition}
\end{definitionbox}

\begin{remark*}[Standard quantified statement]
\[
a := b \quad\text{introduces notation, while}\quad a=b \quad\text{asserts equality.}
\]
\end{remark*}

\begin{remark*}[Interpretation]
The distinction prevents definitions from being confused with theorems.
Definitions create abbreviations; propositional equalities must be available
as axioms, assumptions, or proved statements.
\end{remark*}

\NoLocalDependencies
